\Subhead{1.9.1}{Собственные колебания. Собственные функции системы (моды)}
В распределённой системе координата --- функция пространства: смещение струны, напряжение длинной линии или поле резонатора. Граничные условия отбирают собственные формы и частоты. Это бесконечномерное продолжение~\Qref{1.5.1}.
\begin{equation}\label{eq:distributedmodes}
 \rho(\bm r)u_{tt}+\mathcal Lu=f,\qquad
 \mathcal L\phi_n=\omega_n^2\rho\phi_n.
\end{equation}
При $\rho>0$, самосопряжённых операторе и граничных условиях умножение уравнений для разных мод и вычитание даёт
$(\omega_n^2-\omega_m^2)\int\rho\phi_m^*\phi_n\dd V=0$.
Выберем $\int\rho\phi_m^*\phi_n\dd V=\delta_{mn}$. Свободное поле раскладывается по полному набору мод; каждая координата колеблется со своей частотой.
\begin{equation}
 \rho_\ell u_{tt}-\mathcal T u_{xx}=0,\quad u(0)=u(l)=0:
 \quad \phi_n=\sqrt{\frac2{\rho_\ell l}}\sin\frac{n\pi x}l,\quad
 \omega_n=\frac{n\pi}l\sqrt{\frac{\mathcal T}{\rho_\ell}}.
\end{equation}
Здесь $\mathcal T$ --- натяжение, $\rho_\ell$ --- линейная плотность. Аналогично концы длинной линии задают узлы/пучности напряжения и тока. В пособии Биленко и соавторов этот переход разобран в §2.3.1, с.~76--82; телеграфные уравнения см.~\Qref{4.6.1}.

В идеальном электромагнитном резонаторе
\begin{equation}
 \nabla\times\mu^{-1}\nabla\times\bm E_n
 =\omega_n^2\varepsilon\bm E_n,\qquad
 \bm n\times\bm E_n|_{\partial V}=0.
\end{equation}
Для недиспергирующего диэлектрика естественен вес $\varepsilon$; дополнительно соблюдаются уравнения дивергенции. Открытость, потери и дисперсия могут нарушать самосопряжённость: простую положительную нормировку нельзя переносить без проверки. В неограниченной области появляется непрерывный спектр.
